\documentclass{article}
\usepackage{polski}
\begin{document}
\subsection*{Dziedzina Problemu: Przychodnia Medyczna}
\subsection*{Proces biznesowy: Obsługa wizyt}
\begin{itemize}
\item Czynności:
\begin{enumerate}
\item Pracownik wpisuje wizytę
\item System sugeruje wolne terminy
\item Pracownik potwierdza zgodność danych
\item Pracownik rezerwuje pomieszczenie
\item System przechowuje dane o wizycie
\item W dniu wizyty system udostępnia uprawnionym osobom dane wizyty
\end{enumerate}
\end{itemize}
\subsubsection*{Dziedzina problemu System Informatyczny dla Przychodni:}
\begin{itemize}
	\item Wymagania funkcjonalne
		\begin{enumerate}
			\item System musi umożliwić pracownikom rezerwacje wizyty
			\item System powinien pokazywać dostępne terminy wizyt
			\item System powinien pokazywać dostępne pomieszczenia we wskazanym terminie
			\item System powinien pokazywać dane pacjenta
		\end{enumerate}
	\item Wymagania jakościowe
		\begin{enumerate}
			\item System powinien mieć interfejs użytkownika, który użytkownicy ocenią jako przejrzysty i łatwy w obsłudze
			\item System powinien pokazywać aktualne rezerwacje w co najwyżej 1 minutę
		\end{enumerate}
	\item Ograniczenia
		\begin{enumerate}
			\item
		\end{enumerate}
\end{itemize}



\end{document}